\section{Version history}\label{app:version-history}
\begingroup\small
Every numbered statement of earlier versions keeps its number and
place.  Several statements were renamed in Version~4, as recorded
below.

\subsection*{Version 4 (October 1, 2026)}
\begin{itemize}
    \item The paper was retitled; Versions~1--3 were titled \emph{One
    Meta-Theory, Three Clay-Problem Closures}.  The abstract, the
    introduction and the scope section were rewritten to state the
    status of every result, and a glossary was added.
  \item In \Cref{sec:foreclosure}, targets and accepted imports are now
    claims with an interpretation, instead of propositions.  With
    propositions, a sound set of accepted imports could only accept all
    true propositions or none, and the hypotheses of
    \Cref{cor:bridge-impossibility} could not hold together; the formal
    development now exhibits a model in which they do.
    \Cref{thm:carrier-dichotomy} no longer assumes the aiming, typing
    and coverage-context hypotheses, which its proof does not use.  The
    classification hypothesis (Cls) and the bridge-obligation hypothesis
    (BO), previously applied as unproved obligations, are now explicit
    hypotheses of \Cref{thm:carrier-dichotomy,cor:bridge-impossibility},
        and \Cref{thm:type-finite-foreclosure} assumes coverage explicitly.
    It is shown that the Coverage schema follows from (Cls) when some
    claim expresses the residual closure.  The names of these three
    results now end in ``conditional form''.
  \item The formal record no longer states any schema as an axiom; each
    schema is a definition, used only as a hypothesis.
  \item \Cref{thm:attack-foreclosure-v4,thm:attack-foreclosure-v5},
    previously called Attack Foreclosure v4 and v5, are now called
    Attack Foreclosure and Refined Attack Foreclosure.  Readings of them
    that do not follow from their statements are now described as
    interpretations, and the first clause of
    \Cref{thm:attack-foreclosure-v5} is derived from
    \Cref{thm:attack-foreclosure-v4}.
    \item Definitions~\ref{def:vdiff-tso-column}, \ref{def:applies} and
    \ref{def:load-bearing}, previously called V-Differential
    trace-state-only column, Template-applies predicate and
    Load-bearing predicate, are now called Columns, Template applies and
    Load-bearing content; \Cref{thm:concrete-triad-gaussian-xi},
    previously the Concrete Gaussian--XI triad span, is now the Concrete
        Gaussian--residual triad.  In \Cref{sec:calibration},
    Theorems~\ref{thm:calibration-family-structure},
    \ref{thm:bridges-as-theorems}, \ref{thm:coverage-family-geometry},
    \ref{thm:framework-type-self-char} and
    \ref{cat:named-functorial-gates}, previously called
    Calibration-Family Structure, Bridges-as-Theorems, Coverage-Family
    Geometry, Framework-Type Self-Characterization and Named Functorial
    Gates, are now called Seed catalogue size, Bridge kinds, Declared
    coverage table, Recorded type mismatch and Gate catalogue, names
    that describe what they check.  \Cref{def:load-bearing} defines the
    load-bearing content as an assignment, as in the formal record.  The verdicts of
    \Cref{tab:verdict-signature} are described as this paper's reading
    of the two applications rather than as verified instances.
    \item The SB closure assumption is no longer attributed to
    Foundations~I, the result quoted from Foundations~III is stated
    with its scope, and the name of the Selberg source is described as
    an analogy.
  \item The recognition sources are attributed to the applications
    that name them, the formed-closure predicate is identified as a
    definition of this paper, the sweep reports read in
    \Cref{sec:validation} are cited to the first versions of the
    applications, and the Navier--Stokes physical-window return is
    stated with its time-integrability condition and its assumed
    terminal-divergence input.
  \item \Cref{sec:gaussian-bridge,sec:triad-bridge} now define the
        observation, the heat excess and the frame explicitly, in the
    Sobolev-weighted Fourier coordinates in which they are formalized.  The
    reflection is \(J(\rho)=1-\bar\rho\) (previously misprinted as
    \(1-\rho\)).  \Cref{thm:joint-gaussian-return} is stated for
    nonnegative integer weights, and \Cref{thm:concrete-triad-gaussian-xi}
    for the specific three-solution frame and unit weights for which it
    is proved; the earlier statement for an arbitrary frame was too
    general.
  \item In \Cref{sec:calibration} the conclusions of
    \Cref{thm:calibration-family-structure,thm:bridges-as-theorems,thm:bridge-type-direction}
    are restricted to the properties of records that are proved; the
    seed catalogue is described as introduced in this paper, and its
    source column was removed.
    \item The account of the non-descending translation statement, an
    earlier formulation from which the calibration family developed,
    was reduced to a statement of scope in \Cref{sec:scope}.
  \item Companion papers are cited in their current versions, and the
    first versions separately where their content is meant.
  \item A data point on the depth of matrix-element towers, not
    supported by the cited papers, was removed from
    \Cref{sec:predictions}.
    \item \Cref{app:formalization} was rewritten as a correspondence
    table with a statement of what is and is not verified.
    Cross-references now name the kind of statement they point to.
\end{itemize}

\subsection*{Versions 2 and 3}
Versions~2 and~3, the latter dated September 29, 2026, restated the
premises of the three companion applications in their current form,
added the Gaussian constructions of \Cref{sec:gaussian-bridge} and
\Cref{sec:triad-bridge} with their figures, and described their formal
verification.

\subsection*{Version 1 (June 17, 2026)}
The first version presented the foreclosure, recognition-mode and
calibration-family parts.
\par\endgroup
